\paragraph{C.3ac (Known finite-field theorem and proved deductions; the strengthened cycle-class conclusion uses the cited 2026 input).}
Let $k=\overline{\mathbb F}_p$, let $A/k$ be an abelian variety of
dimension $g$, and let $h$ be any ample divisor class. Independently
of the 2026 input in D.4c--e, we have
\begin{equation}\label{eq:C3ac-classical}
 \begin{gathered}
 \mathrm{CH}^r(A)_{\mathbb Q,(s)}=0\quad(s\ne0),\qquad
 J^r(A)=0,\\
 h^{g-2r}:\mathrm{CH}^r(A)_{\mathbb Q}
                 \xrightarrow{\ \sim\ }
                 \mathrm{CH}^{g-r}(A)_{\mathbb Q}
                 \quad(2r\le g).
 \end{gathered}
\end{equation}
Consequently both abelian injectivity assertions of C.3h hold
over this field; the first map is an isomorphism. The letter's
$A$, $A'$ and $A''$ hold for these abelian varieties, with the
same rational coefficients and degree ranges as in C.1.

If we additionally invoke Theorem I of
\citet{D_VanHoftenSempliner2026}, with exactly the source qualification
in D.4c--e, then for every $r$ we also have
\begin{equation}\label{eq:C3ac-stronger}
 \mathrm{CH}^r(A)_{\mathbb Q}
       \xrightarrow{\ \sim\ }
       \mathrm{CH}^r(A)_{\mathrm{num},\mathbb Q},
 \qquad G^r(A)=0.
\end{equation}
Under this same input the letter's strict $B$ holds for $A$.
The analogous rational--numerical comparison holds on every
product of abelian varieties and, more generally, on every
smooth projective variety over a finite field whose \emph{Chow}
motive is of abelian type. In the category generated by these
Chow motives, numerical passage is an equivalence onto the
corresponding semisimple numerical category. This last assertion
does not include correspondences from arbitrary varieties
outside that category.

\begin{proof}
\emph{The unconditional statement.}
Choose a finite-field model $A_0/\mathbb F_q$ of $A$.
\citet[Theorem 7.1, pp.~99--100]{cyKunnemann1993} proves
\[
       \mathrm{CH}^r(h^i(A_0))_{\mathbb Q}=0
                    \qquad(i\ne2r).
\]
The same theorem applies after every finite extension of
$\mathbb F_q$. A cycle and a rational-equivalence witness
over $k$ involve finitely many coefficients, so ordinary Chow
groups commute with this filtered union of finite fields.
The canonical multiplication projectors commute with base
change, and the index convention is $i=2r-s$. This gives the
first assertion in \eqref{eq:C3ac-classical}. The middle
assertion is C.5a: every algebraically trivial class is torsion.

The ample line bundle also descends after a finite extension.
\citet[Theorem 7.3, p.~100]{cyKunnemann1993} proves the
displayed hard Lefschetz isomorphism over each such finite
field. Passing to the same filtered union proves it over $k$.
Theorem 7.2 is \emph{not} this result: the source's item 7.2
is a definition of purity. The original result appears in
\citet[\S3.4, Theorem 7, pp.~339--341]{cySoule1984}.
Since $J^r(A)=0$, the quotient $B^r(A)$ agrees with the
ordinary rational Chow group. This proves both C.3h assertions.

For completeness, let $i:Y\hookrightarrow A$ be any smooth
divisor representing $h$, and take $z\in\mathrm{CH}^{r+1}(A)_{
\mathbb Q}$ with $2r\ge g-1$. Hard Lefschetz writes
\[
 z=h^{2r+2-g}a,\qquad
 a\in\mathrm{CH}^{g-r-1}(A)_{\mathbb Q}.
\]
The exponent $2r+1-g$ is nonnegative, and the projection
formula gives
\[
 z=i_*i^*(h^{2r+1-g}a).
\]
This proves $A'$ in this case. Ordinary Chow localization
then gives $A''$ for $A\setminus Y$, and hence for its
quotient modulo algebraic equivalence as well. There is no
appeal here to a localization sequence for homological
equivalence.

\emph{The stronger conclusion and its additional input.}
Theorem I cited above, together with
\citet[Theorem 7.1, author p.~25]{D_Milne1999}, supplies the
Tate conjecture for every finite-field abelian variety and
every finite-field model of every product. We retain Milne's
cofinal-CM-field correction exactly as in D.4e. Tate
surjectivity passes through direct sums, Tate twists, and
algebraic projectors, so it holds for the stated Chow motives
of abelian type, including Artin factors after finite descent.

Now \citet[Definition 1, p.~978; Theorem 1.10, p.~982]{cyKahn2003}
proves rational equals numerical equivalence on each such
finite-field variety. The theorem requires both abelian type
and Tate; K\"unnemann's grading statement alone supplies
neither the missing grade-zero injectivity nor this conclusion.
The finite-dimensionality and nilpotence inputs used by Kahn
are explicit in his Theorem 1.4 and Proposition 1.5, p.~981.
His Lemma 1.6 and Proposition 1.7 then eliminate all nonunit
simple summands that could contribute a homologically trivial
map from the unit. Thus the application uses his full
rational-equivalence theorem, rather than merely a statement
about numerical and homological equivalence.

Let $z\in\mathrm{CH}^r(A)_{\mathbb Q}$ be numerically trivial.
Represent it over a finite extension $\mathbb F_{q^e}$.
Its intersections with all cycles defined over that finite
field vanish, because the corresponding geometric
intersection numbers vanish. Kahn's theorem makes the
representative zero over that finite field. This proves
the injectivity in \eqref{eq:C3ac-stronger}; surjectivity is
the definition of the numerical quotient. A homologically
trivial class is numerically trivial, so $G^r(A)=0$ follows.
The same argument handles every product. Finite data descend
for Chow projectors and morphisms too, giving full
faithfulness after passing to $k$. Splitting numerical
idempotents now produces the claimed equivalence of the
generated pseudo-abelian categories; numerical semisimplicity
is Jannsen's theorem
\citep[Theorem 1, pp.~448--449]{Jannsen1992}.

Finally, for $g=2m$, the primitive middle space modulo
algebraic equivalence identifies, by
\eqref{eq:C3ac-stronger}, with its space of primitive algebraic
cohomology classes. Under the same Hodge-CM input,
\citet[Theorem 3.3 and Corollary 3.4, pp.~607--608]{D_Milne2002}
give the Hodge standard conjecture for abelian varieties.
Thus $(-1)^m\deg(x^2)$ is strictly positive for every nonzero
primitive class, which is the letter's $B$. This positivity
step is stated separately from the classical hard Lefschetz
statement.
\end{proof}

\emph{Scope of the function-field citation.}
\citet[\S8, pp.~101--102]{cyKunnemann1993} relates the
hypercube assertion to negative Beauville vanishing. The
opened original
\citet[\S3.5.5, Corollary 3, p.~343]{cySoule1984}
has a precise constant-variety hypothesis: its varieties over
the function field are $Y_j=X_j\otimes K$ for varieties
$X_j$ over a finite field. We do not remove that hypothesis
when reading the shorter function-field description in
K\"unnemann's discussion. In particular, the proof here
does not settle C.3h for arbitrary nonconstant abelian
varieties over $\overline{\mathbb F_p(t)}$, nor does its
restricted semisimple category settle C.1n with unrestricted
source varieties.
